\section{Zu dieser Anleitung}
\label{sec:about}

Diese Anleitung beschreibt die Bedienung des \ST{} für Anwender (\autoref{sec:overview} bis
\autoref{sec:specs}) sowie die Servicearbeiten Kalibrierung und Firmware-Update
(\autoref{sec:service}). Sie gilt für die auf dem Titelblatt genannten Hardware- und
Firmwareversionen. Die Firmwareversion eines Geräts wird unter \disp{Settings > About} angezeigt.

\subsection{Darstellung}

\begin{tabular}{ll}
    \btn{OUT}                  & Taste am Gerät \\
    \enc                       & Drehgeber (drehen) \\
    \disp{Voltage Source}      & Text auf dem Display \\
    \tbc{60}                   & Wert \TbcName{} (am Gerät noch nicht verifiziert) \\
\end{tabular}

\begin{note}
    Die Bedienoberfläche des Geräts ist englisch. Displaytexte werden deshalb im Original wiedergegeben.
\end{note}

\subsection{Sicherheitshinweise}

\begin{danger}
    Unmittelbare Lebensgefahr. Missachtung führt zu Tod oder schweren Verletzungen.
\end{danger}
\begin{warning}
    Mögliche Lebensgefahr. Missachtung kann zu Tod oder schweren Verletzungen führen.
\end{warning}
\begin{caution}
    Missachtung kann zu leichten Verletzungen oder Schäden am Gerät oder am Prüfling führen.
\end{caution}
\begin{note}
    Nützliche Information zur Bedienung.
\end{note}

\subsection{Abkürzungen}
\begin{acronym}[ESC]
    \acro{BMS}{Batteriemanagementsystem}
    \acro{DUT}{Prüfling (device under test)}
    \acro{FS}{Formula Student}
    \acro{HV}{Hochspannung}
    \acro{LiFePO4}[LiFePO$_4$]{Lithium-Eisenphosphat}
    \acro{OVP}{Überspannungsschutz}
    \acro{SoC}{Ladezustand (state of charge)}
\end{acronym}

\subsection{Änderungsverzeichnis}
\begin{tabular}{lll}
    \rowcolor{gray} \textbf{Revision} & \textbf{Datum} & \textbf{Änderung} \\
    0.1 & 2026-10-24 & Erster Entwurf (Struktur, Inhalt aus Firmware \texttt{sw\_rework}) \\
\end{tabular}
